\section{Introducción}
El problema de las $n$ reinas pregunta cuántas maneras hay de colocar $n$ reinas en un tablero de
$n\times n$ sin que dos se ataquen. Denotamos ese número por $Q(n)$ y contamos por separado las
soluciones relacionadas por una simetría. Numeramos filas y columnas desde uno, con filas de arriba
abajo y columnas de izquierda a derecha, y $(i,j)$ es la casilla de la fila $i$ y la columna $j$.
Leer las columnas de las reinas por orden de filas identifica cada solución con una permutación
$\sigma$, cuyas reinas ocupan las casillas $(i,\sigma(i))$. Las restricciones diagonales piden que
las diferencias $i-\sigma(i)$ sean distintas entre sí, y también las sumas $i+\sigma(i)$; así
determinan cuáles de las $n!$ permutaciones sirven. Por ejemplo, $(2,4,1,3)$ y $(3,1,4,2)$ son las
dos soluciones del tablero de $4\times4$.

Para estudiar cómo crece $Q(n)$, conviene tomar logaritmos. Entre los antecedentes combinatorios
figuran la cota superior de Luria~\cite{luria}, la inferior de Luria y Simkin~\cite{ls} y el
trabajo de Bowtell y Keevash~\cite{bk}; juntos resuelven una conjetura de Rivin, Vardi y
Zimmermann~\cite{rvz} sobre el crecimiento de $Q(n)$. Bell y Stevens~\cite{bs} dan una visión de conjunto.
La fórmula variacional de Simkin~\cite{simkin} determina $\log Q(n)$ hasta un error $o(n)$.
En la normalización del perfil variacional~\cite{pv}, el término principal es $\log(n!)+n\Phi$.
Simkin escribe $Q(n)=((1+o(1))ne^{-\alpha})^n$ con $1{,}94<\alpha<1{,}9449$; como
$\log(n!)=n\log n-n+O(\log n)$, las dos normalizaciones se relacionan por $\Phi=1-\alpha$.
Nobel, Agrawal y Boyd~\cite{nab} obtuvieron cotas numéricas para la constante de primer orden.
Más allá del intervalo enumerado, se han obtenido estimaciones numéricas del crecimiento de $Q(n)$ por
conteo de Monte Carlo~\cite{zm} y, recientemente, por integración termodinámica de un modelo de gas
reticular~\cite{llw}.
El perfil describe la distribución de masa en el tablero y sobre cada diagonal. Estos datos
entran en el predictor canónico $Q_{\rm pred}(n)$ de~\cite{pc}. A partir de los pesos de las casillas
se calcula una entropía y se añade una corrección por la dependencia entre las líneas. El
Apéndice~\ref{app:predictor} da la fórmula finita que usamos en todo el artículo.

El siguiente avance fue identificar el término constante. Después de restar $\log(n!)+n\Phi$,
el logaritmo del predictor tiende a $C_{\rm pred}$~\cite{pc}. El conteo exacto tiene esa misma
constante~\cite{tc}. Por tanto, al formar la diferencia con signo
\[
 R_Q(n)=\log Q(n)-\log Q_{\rm pred}(n),
\]
se cancelan tanto el término principal como la constante. Un paso previo del presente trabajo,
mediante estimaciones de magnitud, da $|R_Q(n)|\le C/n$ para $n$ suficientemente grande.
El Lema~\ref{lem:previa} prueba esa cota a un solo corte fijo, bajo la ley adaptativa, que revela las
filas en orden aleatorio y elige cada reina con probabilidad proporcional al predictor del tablero
resultante (Sección~\ref{sec:signedmap}).
El Apéndice~\ref{app:signed} la usa al desarrollar la identidad de historias detenidas.
Así, $nR_Q(n)$ está acotado, aunque todavía podría oscilar con el tamaño del tablero. ¿Converge?
El Teorema~\ref{teo:principal} identifica ese límite, que denotamos por \(\cerr\). Mientras que
\(C_{\rm pred}\) es el término constante que comparten los dos logaritmos, \(\cerr/n\) es la primera
corrección relativa al pasar del predictor finito al conteo exacto. Si se quisiera desarrollar
directamente \(\log Q(n)\) respecto de \(\log(n!)+n\Phi+C_{\rm pred}\), habría que sumar a
\(\cerr\) el coeficiente de \(1/n\) del propio \(\log Q_{\rm pred}(n)\), si ese coeficiente existe.

\paragraph{En qué se apoya la demostración.} La demostración del Teorema~\ref{teo:principal} no es
autocontenida. Usa la existencia y la regularidad del perfil variacional~\cite{pv}, el corte temprano
de~\cite{ef}, los teoremas de inserción y de consistencia del predictor~\cite{pc}, y la regla de parada,
la permanencia y las cotas estáticas en que se basa el término constante~\cite{tc}. Estos cuatro
trabajos son preprints del autor; cada registro de Zenodo contiene el original en español y una versión
en inglés. La Sección~\ref{sec:importados} enumera los enunciados que se usan, con sus hipótesis, e
indica los dos lugares en que se toma un argumento de una demostración y no sólo de un enunciado. Dos
pasos están asistidos por ordenador, en aritmética racional exacta (Apéndice~\ref{app:pair}). El valor
numérico de \(\cerr\) es una estimación y no forma parte del teorema.

La dificultad se ve al revelar una solución fila por fila. Mientras quedan muchas filas, cada
inserción introduce un defecto pequeño y es su suma con signo la que importa. Al quedar pocas
filas, dos casillas aún disponibles pueden compartir una diagonal. Los resultados de completabilidad
de~\cite{ef} describen el papel de las últimas filas;
Glock, Munhá Correia y Sudakov~\cite{gms} estudiaron cuándo puede completarse una colocación
parcial. Para calcular el coeficiente necesitamos además el peso de esos conflictos en el error.
Siguiendo el corte temprano de~\cite{ef}, detenemos el proceso cuando queda una fracción fija
de filas vacías. Esto permite estudiar el defecto acumulado antes del corte y medir el error en
el tablero residual. La comparación de dos cortes reunirá ambas contribuciones.
La Figura~\ref{fig:mascara-ilustrada} sigue la primera solución de cuatro reinas: al revelar una
sola reina ya se ve la diferencia entre una completación de torres y una de reinas.
\begin{figure}[htbp]
\centering
\begin{tikzpicture}[x=.61cm,y=.61cm]
  \bcuadricula{4}
  \foreach \r/\c in {1/2,2/4,3/1,4/3}{\breina{4}{\r}{\c}}
  \node[font=\small] at (2,-.62) {Solución completa};
  \begin{scope}[shift={(7,0)}]
    \bcuadricula{4}
    \fill[orange,opacity=.24] (0,3) rectangle (4,4);
    \fill[blue,opacity=.15] (1,0) rectangle (2,4);
    \begin{scope}
      \clip (0,0) rectangle (4,4);
      \draw[teal!75!black,line width=.7pt] (-2,7)--(6,-1);
      \draw[red!70!black,dashed,line width=.7pt] (-2,0)--(6,8);
    \end{scope}
    \foreach \r/\c in {2/4,3/1,3/3,4/1,4/3,4/4}{\bmarca{4}{\r}{\c}}
    \breina{4}{1}{2}
    \node[font=\small] at (2,-.62) {Una reina revelada};
  \end{scope}
  \begin{scope}[shift={(14,.1)}]
    \foreach \r/\cells in {1/{0,0,1},2/{1,1,0},3/{1,1,1}}{
      \foreach \w [count=\c] in \cells {
        \ifnum\w=1 \fill[teal!16] (\c-1,3-\r) rectangle (\c,4-\r);\fi
        \node[font=\small] at (\c-.5,3.5-\r) {$\w$};
      }
    }
    \draw[step=1,black!45,line width=.3pt] (0,0) grid (3,3);
    \draw[btinta,line width=.6pt] (0,0) rectangle (3,3);
    \foreach \i/\r in {1/2,2/3,3/4}{\node[font=\scriptsize] at (-.3,3.5-\i) {$\r$};}
    \foreach \i/\c in {1/1,2/3,3/4}{\node[font=\scriptsize] at (\i-.5,3.32) {$\c$};}
    \node[font=\small] at (1.5,-.72) {Máscara residual};
  \end{scope}
\end{tikzpicture}
\caption{De una solución al problema de torres. Tras revelar $(1,2)$, los puntos del centro
son los unos de la máscara de la derecha, con sus etiquetas originales. La fila 2 exige la
columna 4. Las filas 3 y 4 pueden repartirse las columnas 1 y 3 de dos maneras:
$\operatorname{per}M=2$. Sólo la del panel izquierdo completa las reinas. En la otra,
$(2,4)$ y $(3,3)$ se atacan. Así, $Q(B)=1$.}\label{fig:mascara-ilustrada}
\end{figure}

